
\paragraph{First order distribution function}

The initial distribution function for the first order drift kinetic equation
is determined in the non-relativistic limit using the simplified Lorentz
collision model corresponding to $Z_{i}\gg 1$ and $T_{i}=0$. In that case,
the bounce-averaged collision operator reduces to 
\begin{equation}
\left\{ \nabla _{\mathbf{p}}\cdot \mathbf{S}_{\mathbf{p}}\left( f_{1}\right)
\right\} _{\mathcal{L}}=-\frac{p}{\lambda }\frac{\partial }{\partial \xi _{0}%
}\left( \sqrt{1-\xi _{0}^{2}}\lambda \widetilde{S}_{\xi _{0},\mathcal{L}%
}^{\left( 0\right) }\right) -\frac{p}{\lambda }\frac{\partial }{\partial \xi
_{0}}\left( \sqrt{1-\xi _{0}^{2}}\lambda S_{\xi _{0},\mathcal{L}}^{\left(
0\right) }\right)
\end{equation}
and in the steady-state regime $\left\{ \nabla _{\mathbf{p}}\cdot \mathbf{S}%
_{\mathbf{p}}\left( f_{1}\right) \right\} _{\mathcal{L}}=0$ so that the
equation to be solved is 
\begin{equation}
\frac{\partial }{\partial \xi _{0}}\left( \sqrt{1-\xi _{0}^{2}}\lambda
S_{\xi _{0},\mathcal{L}}^{\left( 0\right) }\right) =-\frac{\partial }{%
\partial \xi _{0}}\left( \sqrt{1-\xi _{0}^{2}}\lambda \widetilde{S}_{\xi
_{0},\mathcal{L}}^{\left( 0\right) }\right)
\end{equation}
where 
\begin{equation}
\widetilde{S}_{\xi _{0},\mathcal{L}}^{\left( 0\right) }=\sigma \left\{ \frac{%
\sigma \xi }{\sqrt{\Psi }\xi _{0}}S_{\xi ,\mathcal{L}}\left( \widetilde{f}%
\right) \right\}
\end{equation}
and 
\begin{equation}
S_{\xi _{0},\mathcal{L}}^{\left( 0\right) }=\sigma \left\{ \frac{\sigma \xi 
}{\sqrt{\Psi }\xi _{0}}S_{\xi ,\mathcal{L}}\left( g\right) \right\}
\end{equation}
as defined in Sec.***.

Since only pitch-angle scattering takes place in this limit, $\widetilde{S}%
_{\xi _{0},\mathcal{L}}^{(0)}$ may be expressed in a very simple form 
\begin{equation}
\widetilde{S}_{\xi _{0},\mathcal{L}}^{(0)}=\frac{\sqrt{1-\xi _{0}^{2}}}{p}%
\widetilde{D}_{\xi \xi ,\mathcal{L}}^{\text{C}(0)}\frac{\partial \widetilde{f%
}^{(0)}}{\partial \xi _{0}}
\end{equation}
since $F_{\xi }=D_{\xi p}=0$ for collisions with 
\begin{equation}
\widetilde{D}_{\xi \xi ,\mathcal{L}}^{\text{C}(0)}=\sigma \left\{ \frac{%
\sigma \xi ^{3}}{\Psi ^{2}\xi _{0}^{3}}D_{\xi \xi ,\mathcal{L}}^{\text{C}%
}\right\}
\end{equation}
while 
\begin{equation}
S_{\xi _{0},\mathcal{L}}^{(0)}=\frac{\sqrt{1-\xi _{0}^{2}}}{p}D_{\xi \xi ,%
\mathcal{L}}^{\text{C}(0)}\frac{\partial g^{(0)}}{\partial \xi _{0}}
\end{equation}
with 
\begin{equation}
D_{\xi \xi ,\mathcal{L}}^{\text{C}(0)}=\left\{ \frac{\xi ^{2}}{\Psi \xi
_{0}^{2}}D_{\xi \xi ,\mathcal{L}}^{\text{C}}\right\}
\end{equation}
where $D_{\xi \xi ,\mathcal{L}}^{\text{C}}=B_{t}^{\mathcal{L}}$ according
usual notation given in Sec.***.

The differential equation to be solved is then 
\begin{equation}
\frac{\partial }{\partial \xi _{0}}\left( \left( 1-\xi _{0}^{2}\right)
\lambda D_{\xi \xi ,\mathcal{L}}^{\text{C}(0)}\frac{\partial g^{(0)}}{%
\partial \xi _{0}}\right) =-\frac{\partial }{\partial \xi _{0}}\left( \left(
1-\xi _{0}^{2}\right) \lambda \widetilde{D}_{\xi \xi ,\mathcal{L}}^{\text{C}%
(0)}\frac{\partial \widetilde{f}^{(0)}}{\partial \xi _{0}}\right)
\end{equation}
starting from the known expression of $\partial \widetilde{f}^{(0)}/\partial
\xi _{0}.$ In the Lorentz limit, since $B_{t}^{\mathcal{L}}$ is simply set
to $1/2$, expressions of $D_{\xi \xi ,\mathcal{L}}^{\text{C}(0)}$ and $%
\widetilde{D}_{\xi \xi ,\mathcal{L}}^{\text{C}(0)}$ are obtained in a
straightforward manner, 
\begin{equation}
D_{\xi \xi ,\mathcal{L}}^{\text{C}(0)}=\frac{1}{2}\left\{ \frac{\xi ^{2}}{%
\Psi \xi _{0}^{2}}\right\} =\frac{1}{2}\frac{\lambda _{2,-1,0}}{\lambda }
\end{equation}
and 
\begin{equation}
\widetilde{D}_{\xi \xi ,\mathcal{L}}^{\text{C}(0)}=\frac{1}{2}\sigma \left\{ 
\frac{\sigma \xi ^{3}}{\Psi ^{2}\xi _{0}^{3}}\right\} =\frac{1}{2}\frac{%
\overline{\lambda }_{3,-2,0}}{\lambda }
\end{equation}
and $g^{(0)}$ is given by the simple equation 
\begin{equation}
\frac{\partial }{\partial \xi _{0}}\left( \left( 1-\xi _{0}^{2}\right)
\lambda _{2,-1,0}\frac{\partial g^{(0)}}{\partial \xi _{0}}\right) =-\frac{%
\partial }{\partial \xi _{0}}\left( \left( 1-\xi _{0}^{2}\right) \overline{%
\lambda }_{3,-2,0}\frac{\partial \widetilde{f}^{(0)}}{\partial \xi _{0}}%
\right)
\end{equation}

If the solution of the zero-order Fokker-Planck equation is 
\begin{equation}
f_{0}^{\left( 0\right) }=f_{M}\simeq \frac{n_{e}}{\left[ 2\pi T_{e}\right]
^{3/2}}\exp \left[ -\frac{p^{2}}{2T_{e}}\right]
\end{equation}
its spatial derivative is 
\begin{equation}
\frac{\partial f_{0}^{\left( 0\right) }}{\partial \psi }=\left[ \frac{d\ln
n_{e}}{d\psi }+\left( \frac{p^{2}}{2T_{e}}-\frac{3}{2}\right) \frac{d\ln
T_{e}}{d\psi }\right] f_{M}
\end{equation}
and from the definition of $\widetilde{f}^{(0)}$ given in Sec.***, 
\begin{eqnarray}
\widetilde{f}^{(0)} &=&\frac{p\xi _{0}I\left( \psi \right) }{q_{e}B_{\min
}\left( \psi \right) }\frac{\partial f_{0}^{(0)}(p,\xi _{0},\psi )}{\partial
\psi }  \nonumber \\
&=&\frac{p\xi _{0}I\left( \psi \right) }{q_{e}B_{\min }\left( \psi \right) }%
\left[ \frac{d\ln n_{e}}{d\psi }+\left( \frac{p^{2}}{2T_{e}}-\frac{3}{2}%
\right) \frac{d\ln T_{e}}{d\psi }\right] f_{M}
\end{eqnarray}

Projected on the numerical grids, the discrete form of $\widetilde{f}^{(0)}$
is 
\begin{eqnarray}
\widetilde{f}_{l+1/2,i+1/2,j+1/2}^{(0)} &=&\frac{p_{i+1/2}\xi
_{0,j+1/2}I_{l+1/2}}{q_{e}B_{\min ,l+1/2}}\left[ \left. \frac{d\ln n_{e}}{%
d\psi }\right| _{l+1/2}\right.  \nonumber \\
&&\left. +\left( \frac{p_{i+1/2}^{2}}{2T_{e,l+1/2}}-\frac{3}{2}\right)
\left. \frac{d\ln T_{e}}{d\psi }\right| _{l+1/2}\right]
f_{M,l+1/2,i+1/2,j+1/2}
\end{eqnarray}
where radial derivatives are performed as discussed in Sec.***.

Since $f_{M}$ is by definition independent of $\xi _{0},$ 
\begin{eqnarray}
\frac{\partial \widetilde{f}^{(0)}}{\partial \xi _{0}} &=&\frac{pI\left(
\psi \right) }{q_{e}B_{\min }\left( \psi \right) }\left[ \frac{d\ln n_{e}}{%
d\psi }+\left( \frac{p^{2}}{2T_{e}}-\frac{3}{2}\right) \frac{d\ln T_{e}}{%
d\psi }\right] f_{M}  \nonumber \\
&=&\Xi \left( \psi ,p\right) f_{M}
\end{eqnarray}
where 
\begin{equation}
\Xi \left( \psi ,p\right) =\frac{pI\left( \psi \right) }{q_{e}B_{\min
}\left( \psi \right) }\left[ \frac{d\ln n_{e}}{d\psi }+\left( \frac{p^{2}}{%
2T_{e}}-\frac{3}{2}\right) \frac{d\ln T_{e}}{d\psi }\right]
\end{equation}
the equation becomes 
\begin{equation}
\left( 1-\xi _{0}^{2}\right) \lambda _{2,-1,0}\frac{\partial g^{(0)}}{%
\partial \xi _{0}}=-\Xi \left( \psi ,p\right) f_{M}\left( 1-\xi
_{0}^{2}\right) \overline{\lambda }_{3,-2,0}+C
\end{equation}

Assuming $\partial g^{(0)}/\partial \xi _{0}$ is finite at $\left| \xi
_{0}\right| =1,$ $C=0,$ and 
\begin{eqnarray}
\frac{\partial g^{(0)}}{\partial \xi _{0}} &=&-\Xi \left( \psi ,p\right)
f_{M}\frac{\overline{\lambda }_{3,-2,0}}{\lambda _{2,-1,0}}  \nonumber \\
&=&-\Xi \left( \psi ,p\right) f_{M}\frac{\lambda _{3,-2,0}}{\lambda _{2,-1,0}%
}H\left( \left| \xi _{0}\right| -\xi _{0T}\right)
\end{eqnarray}
using the definition of $\overline{\lambda }_{n,m,p}$ given in Sec. ***. By
definition, $\overline{\lambda }_{3,-2,0}$ is an even function of $\xi _{0},$
and therefore $g^{(0)}$ must be an odd function $\xi _{0}.$ Since
furthermore $g^{(0)}$ is an even function of $\xi _{0}$, $g^{(0)}$ $=0$ in
the trapped region. Integrating in the region $-$ $\xi _{0}\leq \xi _{0T},$ 
\begin{eqnarray}
g^{(0)}\left( \psi ,p,\xi _{0}\right) &=&-\Xi \left( \psi ,p\right)
f_{M}\int_{-\xi _{0T}}^{\xi _{0}}\frac{\overline{\lambda }_{3,-2,0}}{\lambda
_{2,-1,0}}d\xi _{0}^{\prime }  \nonumber \\
&=&\Xi \left( \psi ,p\right) f_{M}\int_{\xi _{0T}}^{-\xi _{0}}\frac{%
\overline{\lambda }_{3,-2,0}}{\lambda _{2,-1,0}}d\xi _{0}^{\prime \prime }
\end{eqnarray}
using the relation $\xi _{0}^{\prime \prime }=-\xi _{0}^{\prime }.$
Performing the same operation in the region $\xi _{0}\geq \xi _{0T},$ one
obtains 
\begin{equation}
g^{(0)}\left( \psi ,p,\xi _{0}\right) =-\Xi \left( \psi ,p\right)
f_{M}\int_{\xi _{0T}}^{\xi _{0}}\frac{\overline{\lambda }_{3,-2,0}}{\lambda
_{2,-1,0}}d\xi _{0}^{\prime }
\end{equation}
and gathering results 
\begin{eqnarray}
g^{(0)}\left( \psi ,p,\xi _{0}\right) &=&-\sigma H\left( \left| \xi
_{0}\right| -\xi _{0T}\right) \Xi \left( \psi ,p\right) f_{M}\int_{\xi
_{0T}}^{\left| \xi _{0}\right| }\frac{\overline{\lambda }_{3,-2,0}}{\lambda
_{2,-1,0}}d\xi _{0}^{\prime }  \nonumber \\
&=&-\sigma H\left( \left| \xi _{0}\right| -\xi _{0T}\right) \Xi \left( \psi
,p\right) I_{\mathcal{L}}\left( \psi ,\left| \xi _{0}\right| \right) f_{M}
\end{eqnarray}
where 
\begin{equation}
I_{\mathcal{L}}\left( \psi ,\left| \xi _{0}\right| \right) =\int_{\xi
_{0T}}^{\left| \xi _{0}\right| }\frac{\overline{\lambda }_{3,-2,0}\left(
\psi ,\xi _{0}^{\prime }\right) }{\lambda _{2,-1,0}\left( \psi ,\xi
_{0}^{\prime }\right) }d\xi _{0}^{\prime }
\end{equation}

The discrete form of $g^{(0)}\left( \psi ,p,\xi _{0}\right) $ is 
\begin{equation}
g_{l+1/2,i+1/2,j+1/2}^{(0)}=-\sigma _{j+1/2}H\left( \left| \xi
_{0,j+1/2}\right| -\xi _{0T,l+1/2}\right) \Xi _{l+1/2,i+1/2}I_{\mathcal{L}%
}\left( \psi _{l+1/2},\left| \xi _{0,j+1/2}\right| \right)
f_{M,l+1/2,i+1/2,j+1/2}
\end{equation}
with 
\begin{equation}
I_{\mathcal{L}}\left( \psi _{l+1/2},\left| \xi _{0,j+1/2}\right| \right)
=\sum_{j^{\prime }=j_{T,l+1/2}^{+}-n_{\xi _{0T,l+1/2}}}^{j^{\prime
}=j-n_{\xi _{0T,l+1/2}}}\frac{\overline{\lambda }_{3,-2,0}^{l+1/2,j^{\prime
}+1/2}}{\lambda _{2,-1,0}^{l+1/2,j^{\prime }+1/2}}\Delta \xi _{0,j^{\prime
}+1/2}
\end{equation}
and $j>j_{T,l+1/2}^{+}.$

\paragraph{Case of circular concentric flux surfaces}

According to Sec.***, 
\begin{equation}
\overline{\lambda }_{3,-2,0}=\frac{H(|\xi _{0}|-\xi _{0T})}{\left(
1+\epsilon \right) \xi _{0}^{2}}\left( \xi _{0}^{2}-\frac{\epsilon \left(
1-\epsilon /2\right) }{\left( 1+\epsilon \right) }\right)
\end{equation}
and 
\begin{equation}
\lambda _{2,-1,0}=\lambda \left( 1-\frac{\Delta _{b}}{\xi _{0}^{2}}\right)
\end{equation}

The integral $I_{\mathcal{L}}\left( \psi ,\xi _{0}\right) $ then becomes 
\begin{equation}
I_{\mathcal{L}}\left( \psi ,\left| \xi _{0}\right| \right) =\int_{\xi
_{0T}}^{\left| \xi _{0}\right| }\frac{\left( \xi _{0}^{\prime 2}-\epsilon
\left( 1-\epsilon /2\right) /\left( 1+\epsilon \right) \right) }{\lambda
\left( 1-\Delta _{b}/\xi _{0}^{\prime 2}\right) \left( 1+\epsilon \right)
\xi _{0}^{\prime 2}}d\xi _{0}^{\prime }
\end{equation}
and in the limit $\epsilon \ll 1,$%
\begin{equation}
I_{\mathcal{L}}\left( \psi ,\left| \xi _{0}\right| \right) \simeq \frac{1}{%
1+\epsilon }\int_{\xi _{0T}}^{\left| \xi _{0}\right| }\frac{d\xi
_{0}^{\prime }}{\lambda \left( 1-\Delta _{b}/\xi _{0}^{\prime 2}\right) }
\end{equation}

\paragraph{Flux surface averaged bootstrap current\label{BootLor}}

The flux-averaged electron bootstrap current in the Lorentz model is given
by the relation 
\begin{equation}
\left\langle J_{\Vert ,\mathcal{L}}\right\rangle \left( \psi \right)
=\left\langle \widetilde{J}_{\Vert ,\mathcal{L}}\right\rangle _{\phi
}^{1}\left( \psi \right) +\left\langle J_{\Vert ,\mathcal{L}}\right\rangle
_{\phi }^{1}\left( \psi \right)
\end{equation}
where 
\begin{equation}
\left\langle \widetilde{J}_{\Vert }\right\rangle _{\phi }^{1}\left( \psi
\right) =2\pi q_{e}\frac{\widetilde{q}\left( \psi \right) }{\overline{q}%
\left( \psi \right) }\frac{R_{0}}{R\left( \psi ,\theta _{0}\right) }\frac{%
B_{T}\left( \psi ,\theta _{0}\right) }{B_{\min }\left( \psi \right) }%
\int_{0}^{\infty }dp\;p^{3}\int_{-1}^{1}\;\lambda _{2,-2,2}\xi _{0}%
\widetilde{f}^{(0)}d\xi _{0}
\end{equation}
and

\begin{equation}
\left\langle J_{\Vert }\right\rangle _{\phi }^{1}\left( \psi \right) =2\pi
q_{e}\frac{q\left( \psi \right) }{\overline{q}\left( \psi \right) }%
\int_{0}^{\infty }dp\;p^{3}\int_{-1}^{1}\;H\left( \left| \xi _{0}\right|
-\xi _{0T}\right) \;\xi _{0}g^{(0)}d\xi _{0}
\end{equation}
in the non-relativistic limit $\gamma \rightarrow 1.$ Therefore, 
\begin{eqnarray}
\left\langle \widetilde{J}_{\Vert }\right\rangle _{\phi }^{1}\left( \psi
\right) &=&2\pi \frac{\widetilde{q}\left( \psi \right) }{\overline{q}\left(
\psi \right) }\frac{R_{0}}{R\left( \psi ,\theta _{0}\right) }\frac{%
B_{T}\left( \psi ,\theta _{0}\right) }{B_{\min }^{2}\left( \psi \right) }%
I\left( \psi \right) \times  \nonumber \\
&&\int_{0}^{\infty }\left[ \frac{d\ln n_{e}}{d\psi }+\left( \frac{p^{2}}{%
2T_{e}}-\frac{3}{2}\right) \frac{d\ln T_{e}}{d\psi }\right]
\;p^{4}f_{M}dp\int_{-1}^{1}\xi _{0}^{2}\lambda _{2,-2,2}d\xi _{0}
\end{eqnarray}
and 
\begin{eqnarray}
\left\langle J_{\Vert }\right\rangle _{\phi }^{1}\left( \psi \right)
&=&-2\pi \frac{q\left( \psi \right) }{\overline{q}\left( \psi \right) }\frac{%
I\left( \psi \right) }{B_{\min }\left( \psi \right) }\times \int_{0}^{\infty
}\left[ \frac{d\ln n_{e}}{d\psi }+\left( \frac{p^{2}}{2T_{e}}-\frac{3}{2}%
\right) \frac{d\ln T_{e}}{d\psi }\right] \;p^{4}f_{M}dp  \nonumber \\
&&\;\int_{-1}^{1}\sigma H\left( \left| \xi _{0}\right| -\xi _{0T}\right)
\;\xi _{0}I_{\mathcal{L}}\left( \psi ,\left| \xi _{0}\right| \right) d\xi
_{0}
\end{eqnarray}

Using relation $R\left( \psi ,\theta _{0}\right) B_{T}\left( \psi ,\theta
_{0}\right) =I\left( \psi \right) ,$ one obtains 
\begin{eqnarray}
\left\langle \widetilde{J}_{\Vert }\right\rangle _{\phi }^{1}\left( \psi
\right) &=&2\pi R_{0}\frac{\widetilde{q}\left( \psi \right) }{\overline{q}%
\left( \psi \right) }\frac{B_{T}^{2}\left( \psi ,\theta _{0}\right) }{%
B_{\min }^{2}\left( \psi \right) }I\left( \psi \right) \times  \nonumber \\
&&\int_{0}^{\infty }\left[ \frac{d\ln n_{e}}{d\psi }+\left( \frac{p^{2}}{%
2T_{e}}-\frac{3}{2}\right) \frac{d\ln T_{e}}{d\psi }\right]
\;p^{4}f_{M}dp\int_{-1}^{1}\xi _{0}^{2}\lambda _{2,-2,2}d\xi _{0}
\end{eqnarray}
and 
\begin{eqnarray}
\left\langle J_{\Vert }\right\rangle _{\phi }^{1}\left( \psi \right)
&=&-2\pi \frac{q\left( \psi \right) }{\overline{q}\left( \psi \right) }%
R\left( \psi ,\theta _{0}\right) \frac{B_{T}\left( \psi ,\theta _{0}\right) 
}{B_{\min }\left( \psi \right) }\times \int_{0}^{\infty }\left[ \frac{d\ln
n_{e}}{d\psi }+\left( \frac{p^{2}}{2T_{e}}-\frac{3}{2}\right) \frac{d\ln
T_{e}}{d\psi }\right] \;p^{4}f_{M}dp  \nonumber \\
&&\;\int_{-1}^{1}\sigma H\left( \left| \xi _{0}\right| -\xi _{0T}\right)
\;\xi _{0}I_{\mathcal{L}}\left( \psi ,\left| \xi _{0}\right| \right) d\xi
_{0}
\end{eqnarray}

Since 
\begin{equation}
2\pi \int_{0}^{\infty }\left[ \frac{d\ln n_{e}}{d\psi }+\left( \frac{p^{2}}{%
2T_{e}}-\frac{3}{2}\right) \frac{d\ln T_{e}}{d\psi }\right] \;p^{4}f_{M}dp=%
\frac{3}{2}n_{e}T_{e}\left( \frac{d\ln n_{e}}{d\psi }+\frac{d\ln T_{e}}{%
d\psi }\right)
\end{equation}
using the definite integral relation (Ref.***) 
\begin{equation}
\int_{0}^{\infty }x^{2n}\exp \left( -px^{2}\right) =\frac{\left( 2n-1\right)
!!}{2\left( 2p\right) ^{n}}\sqrt{\frac{\pi }{p}}
\end{equation}
where $\left( 2n-1\right) !!=1.3.5...\times \left( 2n-1\right) ,$ 
\begin{eqnarray}
\left\langle \widetilde{J}_{\Vert }\right\rangle _{\phi }^{1}\left( \psi
\right) &=&\frac{\widetilde{q}\left( \psi \right) }{\overline{q}\left( \psi
\right) }R_{0}\frac{B_{T}^{2}\left( \psi ,\theta _{0}\right) }{B_{\min
}^{2}\left( \psi \right) }\frac{3}{2}n_{e}T_{e}\left( \frac{d\ln n_{e}}{%
d\psi }+\frac{d\ln T_{e}}{d\psi }\right) \times  \nonumber \\
&&\int_{-1}^{1}\xi _{0}^{2}\lambda _{2,-2,2}d\xi _{0}
\end{eqnarray}
and 
\begin{eqnarray}
\left\langle J_{\Vert }\right\rangle _{\phi }^{1}\left( \psi \right) &=&-%
\frac{q\left( \psi \right) }{\overline{q}\left( \psi \right) }R\left( \psi
,\theta _{0}\right) \frac{B_{T}\left( \psi ,\theta _{0}\right) }{B_{\min
}\left( \psi \right) }\frac{3}{2}n_{e}T_{e}\left( \frac{d\ln n_{e}}{d\psi }+%
\frac{d\ln T_{e}}{d\psi }\right) \times  \nonumber \\
&&\;\int_{-1}^{1}\sigma H\left( \left| \xi _{0}\right| -\xi _{0T}\right)
\;\xi _{0}I_{\mathcal{L}}\left( \psi ,\left| \xi _{0}\right| \right) d\xi
_{0}
\end{eqnarray}

Finally, 
\begin{equation}
\left\langle \widetilde{J}_{\Vert ,\mathcal{L}}\right\rangle \left( \psi
\right) =\mathcal{F}_{t}^{eff.}\left( \psi \right) n_{e}T_{e}R_{0}\left( 
\frac{d\ln n_{e}}{d\psi }+\frac{d\ln T_{e}}{d\psi }\right)
\end{equation}
where the flux surface function usualy named ``effective trapped fraction'' $%
\mathcal{F}_{t}^{eff.}\left( \psi \right) $ is 
\begin{eqnarray}
\mathcal{F}_{t}^{eff.}\left( \psi \right) &=&\frac{3}{2}\frac{B_{T}\left(
\psi ,\theta _{0}\right) }{B_{\min }\left( \psi \right) }\times  \nonumber \\
&&\left[ \frac{\widetilde{q}\left( \psi \right) }{\overline{q}\left( \psi
\right) }\frac{B_{T}\left( \psi ,\theta _{0}\right) }{B_{\min }\left( \psi
\right) }\int_{-1}^{1}\xi _{0}^{2}\lambda _{2,-2,2}d\xi _{0}-\frac{q\left(
\psi \right) }{\overline{q}\left( \psi \right) }\frac{R\left( \psi ,\theta
_{0}\right) }{R_{0}}\int_{-1}^{1}\sigma H\left( \left| \xi _{0}\right| -\xi
_{0T}\right) \;\xi _{0}I_{\mathcal{L}}\left( \psi ,\left| \xi _{0}\right|
\right) d\xi _{0}\right]
\end{eqnarray}

In the case of circular concentric flux surfaces, 
\begin{equation}
\frac{R\left( r,\theta _{0}\right) }{R_{0}}=1+\epsilon
\end{equation}
\begin{equation}
\frac{\widetilde{q}\left( r\right) }{\overline{q}\left( r\right) }=1+\epsilon
\end{equation}
and 
\begin{equation}
\frac{q\left( r\right) }{\overline{q}\left( r\right) }=\frac{B_{T}\left(
r,\theta _{0}\right) }{B_{\min }\left( r\right) }\sqrt{\frac{1+\epsilon }{%
1-\epsilon }}
\end{equation}
as shown in Sec. ***. Hence, 
\begin{equation}
\mathcal{F}_{t}^{eff.}\left( \psi \right) =\frac{3}{2}\left( 1+\epsilon
\right) \left( \frac{B_{T}\left( \psi ,\theta _{0}\right) }{B_{\min }\left(
\psi \right) }\right) ^{2}\left[ \int_{-1}^{1}\xi _{0}^{2}\lambda
_{2,0,0}d\xi _{0}-\sqrt{\frac{1+\epsilon }{1-\epsilon }}\int_{-1}^{1}\sigma
H\left( \left| \xi _{0}\right| -\xi _{0T}\right) \;\xi _{0}I_{\mathcal{L}%
}\left( \psi ,\left| \xi _{0}\right| \right) d\xi _{0}\right]
\end{equation}
and since $R\left( \psi ,\theta _{0}\right) /R\left( \psi ,\theta \right)
=\Psi ,$ $\lambda _{2,-2,2}=\lambda _{2,0,0}$, as shown in Sec. ***. In the
limit $B_{p}\ll B_{T},$ further simplifications may be carried out and 
\begin{equation}
\lim_{\epsilon \rightarrow 0}\mathcal{F}_{t}^{eff.}\left( \psi \right)
\simeq \frac{3}{2}\left( 1+\epsilon \right) \left[ \int_{-1}^{1}\xi
_{0}^{2}\lambda _{2,0,0}d\xi _{0}-\sqrt{\frac{1+\epsilon }{1-\epsilon }}%
\int_{-1}^{1}\sigma H\left( \left| \xi _{0}\right| -\xi _{0T}\right) \;\xi
_{0}I_{\mathcal{L}}\left( \psi ,\left| \xi _{0}\right| \right) d\xi
_{0}\right]
\end{equation}

The first term in the square bracket may be evaluated analyticaly. Indeed 
\begin{eqnarray}
\int_{-1}^{1}\xi _{0}^{2}\lambda _{2,0,0}d\xi _{0} &=&\frac{1}{\widetilde{q}}%
\int_{-1}^{1}\xi _{0}^{2}d\xi _{0}\left[ \frac{1}{2}\sum_{\sigma }\right]
_{T}\int_{0}^{2\pi }\frac{d\theta }{2\pi }\epsilon \frac{B}{B_{P}}\frac{\xi
_{0}}{\xi }\frac{\xi ^{2}}{\xi _{0}^{2}}  \nonumber \\
&=&\int_{-1}^{1}d\xi _{0}\left[ \frac{1}{2}\sum_{\sigma }\right]
_{T}\int_{0}^{2\pi }\frac{d\theta }{2\pi }\xi _{0}\xi
\end{eqnarray}
since $B/B_{P}$ is only function of $r$ and $\widetilde{q}=\epsilon B/B_{P}.$
The sum over trapped electrons may be dropped, since $\xi _{0}\xi $ is
independent of $\sigma .$ Therefore, 
\begin{eqnarray}
\int_{-1}^{1}\xi _{0}^{2}\lambda _{2,0,0}d\xi _{0} &=&\int_{-1}^{1}d\xi
_{0}\int_{0}^{2\pi }\frac{d\theta }{2\pi }\xi _{0}\xi  \nonumber \\
&=&\int_{0}^{2\pi }\frac{d\theta }{2\pi }\int_{-1}^{1}\xi _{0}\xi H\left(
\left| \xi _{0}\right| -\sqrt{1-\frac{1}{\Psi }}\right) d\xi _{0}
\end{eqnarray}
which means that only particles who reach the position $\theta $ must be
considered. Since 
\begin{eqnarray}
\int_{-1}^{1}\xi _{0}\xi H\left( \left| \xi _{0}\right| -\sqrt{1-\frac{1}{%
\Psi }}\right) d\xi _{0} &=&2\int_{\sqrt{1-\frac{1}{\Psi }}}^{1}\xi _{0}\xi
d\xi _{0}  \nonumber \\
&=&2\int_{\sqrt{1-\frac{1}{\Psi }}}^{1}\xi _{0}\sqrt{1-\Psi \left( 1-\xi
_{0}^{2}\right) }d\xi _{0}  \nonumber \\
&=&\frac{2}{3}\frac{1}{\Psi }
\end{eqnarray}
and 
\begin{eqnarray}
\int_{0}^{2\pi }\frac{d\theta }{2\pi }\frac{2}{3}\frac{1}{\Psi } &=&\frac{2}{%
3}\int_{0}^{2\pi }\frac{d\theta }{2\pi }\frac{2}{3}\frac{1+\epsilon \cos
\theta }{1+\epsilon }  \nonumber \\
&=&\frac{2}{3}\frac{1}{1+\epsilon }
\end{eqnarray}

Consequently, 
\begin{equation}
\int_{-1}^{1}\xi _{0}^{2}\lambda _{2,0,0}d\xi _{0}=\frac{2}{3}\frac{1}{%
1+\epsilon }
\end{equation}
and using the 
\begin{eqnarray}
\lim_{\epsilon \rightarrow 0}\mathcal{F}_{t}^{eff.}\left( \psi \right)
&\simeq &1-\frac{3}{2}\sqrt{\frac{1}{1-\epsilon ^{2}}}\int_{-1}^{1}\sigma
H\left( \left| \xi _{0}\right| -\xi _{0T}\right) \;\xi _{0}I_{\mathcal{L}%
}\left( \psi ,\left| \xi _{0}\right| \right) d\xi _{0}  \nonumber \\
&\simeq &1-\frac{3}{2}\int_{-1}^{1}\sigma H\left( \left| \xi _{0}\right|
-\xi _{0T}\right) \;\xi _{0}I_{\mathcal{L}}\left( \psi ,\left| \xi
_{0}\right| \right) d\xi _{0}
\end{eqnarray}

It can be shown in Appendix *** after lengthly calculations, that 
\begin{equation}
\lim_{\epsilon \rightarrow 0}\mathcal{F}_{t}^{eff.}\left( \psi \right)
\simeq K\sqrt{\epsilon }+\mathcal{O}\left( \epsilon \right)
\end{equation}
where 
\begin{equation}
K=\frac{3\sqrt{2}}{2}\left[ 1-\sum_{m}^{\infty }\sum_{i_{1}}^{\infty
}\sum_{i_{2}}^{\infty }...\sum_{i_{m}}^{\infty }\frac{%
C_{i_{1}}C_{i_{2}}...C_{i_{m}}}{2^{i_{1}+i_{2}+...+i_{m}}\left[ 2\left(
i_{1}+i_{2}+...+i_{m}\right) -1\right] }\right]
\end{equation}
with 
\begin{equation}
C_{n}=\frac{\left[ \left( 2n\right) !\right] ^{2}}{\left( 2n-1\right)
2^{3n}\left( n!\right) ^{4}}
\end{equation}

A rather fast convergence is obtained with $m$, and for $m=n=6,$ $K\simeq
1.467,$ the usual value found in the litterature on neoclasical theory.
